\subsection{Post-hoc within-review-status audit}
\label{sec:reviewstrata}
The review-status confound does not by itself prove that the classifier
merely detects review status. We rescored the committed cluster-held-out
test split using the frozen study-07 checkpoints and stratified positive
examples by UniProt record prefix, retaining the same reviewed negative
examples in both comparisons. The PepCNN AUROC on the full test set was
$0.9206$ (1{,}461 positives/2{,}925 negatives). Among the 239 reviewed
\texttt{sp|} positives versus the 2{,}925 reviewed negatives, AUROC was
$0.9643$; among 1{,}222 unreviewed \texttt{tr|} positives versus the same
negatives, $0.9121$. PepGNN scored $0.8703$ and $0.8780$, respectively.
The numbers are in \texttt{results/study17\_review\_strata.json} and can
be regenerated with \texttt{studies/study17\_review\_strata.py}.

The reviewed-only PepCNN result shows discriminative signal is not wholly
explained by the review-status prefix, but it does \emph{not} rescue the
original benchmark. Reviewed positives form a small, selected subgroup;
they may have different sequence families, labels, organisms and length
profiles. The model was trained mostly on unreviewed positives and all
reviewed negatives; class-conditional differences unrelated to AMP
function remain. This post-hoc subgroup is not an independent, prospectively
matched test and its AUROC cannot be generalized to all reviewed AMPs.
The proper reviewed-only retraining and matched evaluation remain undone.
